\section*{Chapitre \ref{ch:hf:on-chain-trading} --- Trading on-chain}

\begin{solution}{exo:hf:on-chain-trading:1}
Le prix du pool, 3\,100~\$, dépasse 3\,000~\$ d'une marge supérieure aux frais~: il faut vendre
\[
\Delta x^\ast=\frac{\sqrt{0.997\times1\,000\times3\,100\,000/3\,000}-1\,000}{0.997}=15.05
\]
jetons, pour un profit de 677~\$ avant le gas.
\end{solution}

\begin{solution}{exo:hf:on-chain-trading:2}
$0.04^2/8\times20\,000\,000=\$4\,000$ par jour, 28\,000~\$ par semaine.
\end{solution}

\begin{solution}{exo:hf:on-chain-trading:3}
$1\,000\,000\times0.05-20=\$49\,980$ avant l'\omterm{def:hf:on-chain-trading:bid}{enchère de bundle} (bundle bid).
\end{solution}

\begin{solution}{exo:hf:on-chain-trading:4}
La valeur de chaque searcher ne diffère de celle des autres que par ses coûts~; avec davantage de concurrents, la deuxième meilleure valeur se
rapproche de la meilleure, et dans une enchère au premier prix le gagnant doit enchérir près d'elle~: il ne conserve que son avantage de coût.
\end{solution}

\begin{solution}{exo:hf:on-chain-trading:5}
La loss-versus-rebalancing (LVR) est fixée par l'ampleur du mouvement du prix de marché entre deux arbitrages, qui dépend de la volatilité et de la
taille du pool~; les frais élargissent seulement la bande à l'intérieur de laquelle l'arbitrage ne rapporte pas, ce qui modifie légèrement la perte.
\end{solution}

\begin{solution}{exo:hf:on-chain-trading:6}
Sa couverture se fait sur une plateforme centralisée (CEX), en dehors de la transaction sur la chaîne~: une jambe peut s'exécuter sans l'autre, ce qui laisse
un risque de prix et un risque d'échec de l'inclusion.
\end{solution}

\begin{solution}{exo:hf:on-chain-trading:7}
La loss-versus-rebalancing atteint 96\,300~\$ sur la semaine (d'après la formule~: 112\,000~\$, soit quatre fois le cas de base)~; le fournisseur de liquidité couvert dégage 294\,000~\$
au lieu de 313\,000~\$, les frais plus élevés liés au volume d'arbitrage plus grand compensant une partie de cet écart.
\end{solution}

\begin{solution}{exo:hf:on-chain-trading:8}
11\,400~\$ est la valeur brute avant le gas et les \omterm{def:hf:on-chain-trading:bid}{enchères de bundle}~; après 7\,100~\$ de gas, le gain net des searchers s'élève à 4\,300~\$, et dans une
enchère concurrentielle, l'essentiel revient aux constructeurs~: le searcher conserve son avantage de coût sur le suivant.
\end{solution}

\begin{solution}{pb:hf:on-chain-trading:1}
\begin{enumerate}
\item Maximiser $g\Delta x\,y/(x+g\Delta x)-p\Delta x$~: $g xy/(x+g\Delta x)^2=p$, donc $\Delta x^\ast=(\sqrt{gxy/p}-x)/g$.
\item Voir la \cref{def:hf:on-chain-trading:atomic}~; l'arbitrage CEX--DEX comporte une jambe hors chaîne.
\item Des bots d'arbitrage paient des frais élevés et optimisent leur latence pour passer devant les utilisateurs~; enchères de gas prioritaires~; risque de consensus lié aux frais d'ordonnancement.
\item Du searcher au constructeur avec une enchère~; du constructeur au proposeur avec un paiement~; le meilleur bloc est proposé.
\item Voir la \cref{def:hf:on-chain-trading:bid}.
\item $a+\frac{n-1}{n}(v-a)$~; 93~\%, 97~\% et 99~\%.
\item Son avantage de coût sur le deuxième meilleur~: 6.7~\$ sur 100~\$ face à un concurrent, 1.0~\$ face à dix-neuf.
\item De la même façon~: de nombreux searchers surveillent les facteurs de santé~; la meilleure enchère remporte la liquidation.
\item Voir la \cref{def:hf:on-chain-trading:hedged}~; $\sigma^2V/8$ par unité de temps.
\item De 25\,800~\$ à 29\,500~\$ par semaine, contre 28\,000~\$ pour la formule.
\item L'impact sur les prix payé par les noise traders, de 20\,000~\$ à 22\,500~\$ par semaine~; les noise traders.
\item Il fixe le volume de bruit, alors que le volume réel baisse avec les frais.
\item Les constructeurs reçoivent 93~\% de la valeur du gagnant avec deux searchers et 99~\% avec vingt~; le fournisseur de liquidité couvert dégage 313\,000~\$ sur 20 millions de dollars en une
semaine, à des frais de 0.3~\%, après 26.000~\$ de loss-versus-rebalancing.
\item Des coûts inférieurs à ceux de ses rivaux~: meilleure latence, inventaire, gas.
\item La perte de la victime (990~\$) finance le brut de l'attaquant (700~\$), dont le constructeur retire l'essentiel.
\item Un mistrial en novembre 2025 dans l'affaire MEV-Boost Peraire-Bueno~; l'application de la loi reste incertaine.
\item La liquidité just-in-time.
\item Les frais issus d'un volume réaliste, à comparer à $\sigma^2V/8$.
\item La valeur de l'impact de sa transaction serait remise aux enchères à son profit.
\item Le constructeur, pour l'essentiel.
\end{enumerate}
\end{solution}

\begin{solution}{iq:hf:on-chain-trading:1}
Le pool est périmé~; le premier arbitrage du bloc suivant le ramène à moins des frais près du nouveau prix~; les searchers se font concurrence, et
le constructeur reçoit l'essentiel de la valeur~; les fournisseurs de liquidité supportent la perte.
\iqlookfor{arbitrage, enchère, loss-versus-rebalancing.}
\end{solution}

\begin{solution}{iq:hf:on-chain-trading:2}
L'arbitrage maintient le pool au prix de marché~; pour un pool à produit constant, la valeur vaut $2\sqrt{kP}$, dont la dérivée seconde en $P$ est
$-\frac12\sqrt k P^{-3/2}$~; le taux de perte vaut $-\frac12\sigma^2P^2V^{\prime\prime}=\frac{\sigma^2}{8}V$.
\iqlookfor{la concavité de la valeur du pool.}
\end{solution}

\begin{solution}{iq:hf:on-chain-trading:3}
Diffuser les transactions en attente et l'état~; simuler les bundles candidats sur le dernier état~; estimer la valeur et les enchères des concurrents~; soumettre à
plusieurs constructeurs avant l'échéance~; suivre l'inclusion.
\iqlookfor{simulation sur l'état, échéance, plusieurs constructeurs.}
\end{solution}

\begin{solution}{iq:hf:on-chain-trading:4}
Une position nue sur la plateforme centralisée~: la déboucler aussitôt, ou soumettre à nouveau la jambe sur chaîne si elle reste rentable~; le système doit traiter
l'échec de l'inclusion comme un événement normal.
\iqlookfor{\omterm{def:hf:index-and-basket-arbitrage:legging}{risque d'exécution par jambes} sur chaîne.}
\end{solution}

\begin{solution}{iq:hf:on-chain-trading:5}
Il faut enchérir en dessous de sa valeur d'un montant inférieur à celui qu'on retiendrait face à des pairs~; si la valeur du concurrent est sûrement plus élevée,
il ne faut enchérir que lorsque la sienne en est proche, ou ne pas enchérir.
\iqlookfor{enchères asymétriques.}
\end{solution}

\begin{solution}{iq:hf:on-chain-trading:6}
Face à d'autres enchérisseurs qui offrent $\beta(v)=a+\frac{n-1}{n}(v-a)$, une enchère $b$ gagne avec la probabilité $\bigl(\frac{b-a}{\frac{n-1}{n}(b_{\max}-a)}\bigr)^{n-1}$
à un facteur d'échelle près~; maximiser $(v-b)(b-a)^{n-1}$ donne $b=a+\frac{n-1}{n}(v-a)$.
\iqlookfor{la condition du premier ordre.}
\end{solution}
